\chapter{量化技术}
\label{chap:quantization}

\section{为什么推理需要量化}
\label{sec:why-inference-needs-quantization}

前面两章分别讨论了 PagedAttention 与 FlashAttention。PagedAttention 解决 KV Cache 的动态内存管理问题，FlashAttention 解决 Attention 计算中的 HBM IO 瓶颈。本章进入推理优化的另一个核心方向：\textbf{量化 Quantization}。

量化的基本思想是：用更低 bit-width 的数值格式表示模型权重、激活或 KV Cache，从而降低显存占用、减少内存带宽压力，并在硬件支持时提升计算吞吐。

在大模型推理中，量化的价值尤其明显。以一个 dense Decoder-only 模型为例，如果参数量为 $N$，使用 FP16/BF16 存储权重，则权重显存约为：

[
M_{\mathrm{weight,fp16}} = 2N\ \mathrm{bytes}.
]

如果使用 INT8 权重量化：

[
M_{\mathrm{weight,int8}} \approx 1N\ \mathrm{bytes}.
]

如果使用 INT4 权重量化：

[
M_{\mathrm{weight,int4}} \approx 0.5N\ \mathrm{bytes}.
]

忽略 scale、zero point、alignment 和 packing metadata，仅从权重本体看，INT8 相比 FP16 约节省一半显存，INT4 约节省四分之三显存。对于 70B 参数模型，FP16/BF16 权重大约需要：

[
70\times 10^9\times 2
=====================

140\ \mathrm{GB}.
]

如果权重压到 INT4，理论权重本体约为：

[
70\times 10^9\times 0.5
=======================

35\ \mathrm{GB}.
]

这会直接改变部署形态：原本需要多张 GPU 才能装下的模型，可能可以用更少 GPU 部署；原本显存主要被权重占满的服务，可以留出更多空间给 KV Cache，从而支持更高并发或更长上下文。

\subsection{权重显存}
\label{subsec:weight-memory}

推理阶段不需要保存梯度和优化器状态，显存主要由以下部分组成：

[
M_{\mathrm{inference}}
======================

M_{\mathrm{weights}}
+
M_{\mathrm{KV}}
+
M_{\mathrm{activations}}
+
M_{\mathrm{workspace}}
+
M_{\mathrm{runtime}}.
]

其中权重显存通常是常驻显存，KV Cache 随请求增长，activation 和 workspace 随 batch、sequence length 和 kernel 实现变化。

在 batch 较小、上下文不太长的场景中，权重可能是主要显存占用；在长上下文和高并发 serving 中，KV Cache 可能成为主要显存占用。量化可以作用于不同对象：

\begin{itemize}
\item \textbf{Weight Quantization}：压缩模型权重，例如 W8A16、W4A16；
\item \textbf{Activation Quantization}：压缩中间激活，例如 W8A8；
\item \textbf{KV Cache Quantization}：压缩推理缓存，例如 FP8 KV、INT8 KV；
\item \textbf{Optimizer Quantization}：训练中压缩优化器状态，本章不重点展开。
\end{itemize}

其中推理部署最常见的是 weight-only quantization，即权重量化到 INT8/INT4，激活仍用 FP16/BF16。它实现难度相对较低，精度风险也比同时量化激活更可控。

\begin{table}[H]
\centering
\small
\caption{常见推理量化对象}
\label{tab:quantization-targets}
\begin{tabular}{p{0.22\textwidth}p{0.28\textwidth}p{0.38\textwidth}}
\toprule
\textbf{量化对象} & \textbf{常见格式} & \textbf{主要收益与风险} \
\midrule
权重
& INT8, INT4, NF4, FP8
& 降低常驻显存和权重读取带宽；风险主要是模型精度下降。 \
\midrule
激活
& INT8, FP8
& 降低 GEMM 输入带宽并启用低精度算力；风险是 activation outlier 造成误差放大。 \
\midrule
KV Cache
& INT8, FP8, INT4
& 降低长上下文和高并发的 cache 显存；风险是 decode 质量和长程依赖下降。 \
\midrule
Logits / Sampling
& 通常保持 FP16/FP32
& 采样对数值稳定敏感，通常不作为主要量化对象。 \
\bottomrule
\end{tabular}
\end{table}

\subsection{显存带宽}
\label{subsec:memory-bandwidth}

量化不仅减少显存容量，也减少显存带宽压力。推理中的 Linear 层通常形如：

[
\mat{Y}=\mat{X}\mat{W}.
]

在 decode 阶段，batch 维和 token 维可能很小，权重读取成本会非常重要。每生成一个 token，模型需要经过所有层，每层读取 QKV projection、output projection、FFN up/gate/down projection 等权重。若权重以 FP16 存储，读取字节数大约是 INT8 的 2 倍、INT4 的 4 倍。

当 kernel 是 memory-bound 时，减少权重读取字节数可以直接提升速度。设某层需要读取权重大小 $M_W$，有效 HBM 带宽为 $\beta$，则权重读取时间下界为：

[
T_{\mathrm{read}}
\ge
\frac{M_W}{\beta}.
]

如果量化后权重字节数降低为原来的 $\rho$，则读取时间下界也降低为：

[
T_{\mathrm{read,quant}}
\ge
\frac{\rho M_W}{\beta}.
]

其中：

[
\rho_{\mathrm{INT8}}\approx \frac{1}{2},
\quad
\rho_{\mathrm{INT4}}\approx \frac{1}{4}.
]

当然，实际速度不一定按这个比例提升，因为还要考虑 dequantization、packing/unpacking、Tensor Core 支持、kernel fusion 和 batch size。量化后的 GEMM 如果需要先把 INT4 解包并反量化到 FP16，再执行普通 GEMM，可能会被解包开销抵消。高性能量化推理必须让 dequantization 与 GEMM 融合。

\subsection{compute-bound 与 memory-bound}
\label{subsec:compute-bound-memory-bound-quantization}

量化是否能加速，取决于当前瓶颈。若推理主要受权重读取或 KV Cache 读取限制，量化往往收益明显；若推理主要受计算 FLOPs 限制，则需要硬件真正支持低精度计算，才能获得吞吐提升。

可以用 roofline 直觉理解。算术强度为：

[
I=\frac{\mathrm{FLOPs}}{\mathrm{Bytes}}.
]

量化减少 bytes，因此提高 $I$。如果原 kernel 处于 memory-bound 区域，提高算术强度可以提升性能；如果原 kernel 已经 compute-bound，仅减少 bytes 未必能提升很多，除非低精度格式也提高 compute throughput。

\begin{figure}[H]
\centering
\begin{tikzpicture}[
font=\small,
axis/.style={-{Latex[length=2.2mm]}, thick},
curve/.style={thick},
point/.style={circle, draw, minimum size=0.16cm}
]

```
\node[font=\bfseries] at (0,3.0) {量化通过减少 Bytes 提高算术强度};

\draw[axis] (-5.1,-1.15) -- (5.2,-1.15) node[right] {Arithmetic Intensity};
\draw[axis] (-5.1,-1.15) -- (-5.1,2.25) node[above] {Performance};

\draw[curve, orange!80!black] (-4.8,-0.95) -- (-1.4,1.55);
\draw[curve, orange!80!black] (-1.4,1.55) -- (4.8,1.55);

\node[point, fill=red!20] (fp16) at (-3.2,0.15) {};
\node[anchor=west] at (-3.0,0.15) {FP16 weights};

\node[point, fill=green!20] (int4) at (-0.8,1.05) {};
\node[anchor=west] at (-0.6,1.05) {INT4 weights};

\draw[->, thick, blue!70] (fp16) -- (int4);

\node[align=left, text width=10.5cm] at (0,-2.2) {
  对 memory-bound 推理，量化减少 HBM 读取字节数，使 kernel 向更高算术强度移动。
  但若 dequantization 开销过高，或 kernel 无法利用低精度算力，速度收益会下降。
};
```

\end{tikzpicture}
\caption{从 roofline 视角理解量化的性能收益}
\label{fig:quantization-roofline}
\end{figure}

\subsection{量化与前几章优化的关系}
\label{subsec:quantization-relation-to-previous-chapters}

量化不是孤立优化。它与前几章技术相互影响：

\begin{itemize}
\item 与 \textbf{PagedAttention} 的关系：KV Cache 量化可以进一步降低 paged block 的显存占用，提高可并发请求数；
\item 与 \textbf{FlashAttention} 的关系：FlashAttention 降低 attention 中间 IO，量化降低输入/权重/KV 字节数，两者可互补；
\item 与 \textbf{Tensor Parallel} 的关系：TP 下权重和 KV head 已经切分，量化进一步降低每 GPU shard 的显存和带宽；
\item 与 \textbf{Serving Scheduler} 的关系：量化降低单请求资源占用，改变 admission control、batch size 和成本模型；
\item 与 \textbf{模型架构} 的关系：RMSNorm、GQA、SwiGLU、MoE、RoPE 等结构会影响量化敏感性。
\end{itemize}

因此，量化不能只看模型文件大小，还要看端到端 serving 指标：

[
\text{TTFT},\quad
\text{TPOT},\quad
\text{tokens/s},\quad
\text{P99 latency},\quad
\text{quality},\quad
\text{cost/token}.
]

\section{数值表示与量化基础}
\label{sec:numerical-representation-and-quantization-basics}

量化的本质是用有限个离散值近似连续或高精度值。对深度学习而言，量化通常包含三个步骤：

\begin{enumerate}
\item 选择低精度表示范围；
\item 将高精度 tensor 映射到低精度整数或浮点值；
\item 在计算时直接使用低精度算子，或反量化回高精度参与计算。
\end{enumerate}

\subsection{FP32、FP16、BF16、FP8、INT8、INT4}
\label{subsec:number-formats}

常见数值格式可以分为浮点和整数两类。

\textbf{FP32} 是传统训练和高精度计算格式，精度和动态范围都较高，但显存和带宽开销大。

\textbf{FP16} 使用 16 bit 浮点，显存减半，Tensor Core 支持好，但指数位较少，动态范围小，训练中可能需要 loss scaling。

\textbf{BF16} 也是 16 bit 浮点，它保留与 FP32 类似的指数位宽，因此动态范围较大，但尾数位较少。大模型训练和推理中 BF16 非常常见。

\textbf{FP8} 是 8 bit 浮点格式，常见有 E4M3、E5M2 等变体。FP8 的优点是比 INT8 更像浮点，有指数动态范围；缺点是硬件和软件栈要求更高，并且需要 scale 管理。

\textbf{INT8} 是 8 bit 整数量化格式，范围通常为：

[
[-128,127]
]

或：

[
[0,255].
]

\textbf{INT4} 是 4 bit 整数量化格式，范围通常为：

[
[-8,7]
]

或：

[
[0,15].
]

INT4 的压缩率高，但误差更大，通常需要 per-group scale、精心校准和高质量量化算法。

\begin{table}[H]
\centering
\small
\caption{常见数值格式对比}
\label{tab:numeric-format-comparison-quant}
\begin{tabular}{p{0.16\textwidth}p{0.16\textwidth}p{0.24\textwidth}p{0.32\textwidth}}
\toprule
\textbf{格式} & \textbf{bit 数} & \textbf{典型用途} & \textbf{特点} \
\midrule
FP32
& 32
& 高精度训练、累积
& 动态范围和精度高，显存/带宽开销大。 \
\midrule
FP16
& 16
& 训练/推理
& Tensor Core 支持好，但动态范围较小。 \
\midrule
BF16
& 16
& 大模型训练/推理
& 动态范围接近 FP32，尾数精度低于 FP16。 \
\midrule
FP8
& 8
& 新硬件上的训练/推理
& 有指数动态范围，需要 scale 与 amax 管理。 \
\midrule
INT8
& 8
& W8A8、weight-only
& 整数表示，硬件支持较广，校准较重要。 \
\midrule
INT4
& 4
& weight-only LLM 推理
& 压缩率高，需要 group scale 和专用 kernel。 \
\bottomrule
\end{tabular}
\end{table}

\subsection{uniform quantization}
\label{subsec:uniform-quantization}

最基础的是 uniform quantization，即用等间隔离散值近似连续值。给定实数 $x$，scale 为 $s$，zero point 为 $z$，量化为整数 $q$：

[
q
=

\operatorname{clip}
\left(
\left\lfloor \frac{x}{s}\right\rceil + z,
q_{\min},
q_{\max}
\right).
]

反量化为：

[
\hat{x}
=======

s(q-z).
]

其中：

\begin{itemize}
\item $s$ 控制量化步长；
\item $z$ 控制实数 0 对应的整数位置；
\item $q_{\min},q_{\max}$ 由 bit-width 决定；
\item $\lfloor\cdot\rceil$ 表示四舍五入。
\end{itemize}

量化误差为：

[
e=x-\hat{x}.
]

Uniform quantization 的优点是简单、硬件友好；缺点是当 tensor 分布有 outlier 时，scale 被 outlier 拉大，大部分普通值的量化分辨率变差。

\subsection{symmetric 与 asymmetric quantization}
\label{subsec:symmetric-asymmetric-quantization}

\textbf{Symmetric quantization} 使用零点 $z=0$，量化区间关于 0 对称。例如 INT8 symmetric 范围为：

[
q\in[-127,127]
]

或：

[
q\in[-128,127].
]

Scale 可取：

[
s
=

\frac{\max(|x_{\min}|,|x_{\max}|)}{q_{\max}}.
]

量化为：

[
q
=

\operatorname{clip}
\left(
\left\lfloor \frac{x}{s}\right\rceil,
q_{\min},
q_{\max}
\right).
]

Symmetric quantization 的优点是计算简单，尤其适合权重，因为权重分布通常近似以 0 为中心。

\textbf{Asymmetric quantization} 使用非零 zero point，使实数范围 $[x_{\min},x_{\max}]$ 映射到整数范围 $[q_{\min},q_{\max}]$：

[
s=
\frac{x_{\max}-x_{\min}}
{q_{\max}-q_{\min}},
]

[
z=
q_{\min}
--------

\left\lfloor
\frac{x_{\min}}{s}
\right\rceil.
]

Asymmetric quantization 更适合非对称分布，例如 ReLU 后激活。但它在矩阵乘法中需要处理 zero point 修正，kernel 更复杂。

\begin{figure}[H]
\centering
\begin{tikzpicture}[
font=\small,
axis/.style={-{Latex[length=2.1mm]}, thick},
tick/.style={thick},
qpt/.style={circle, draw, minimum size=0.18cm, fill=blue!12}
]

```
\node[font=\bfseries] at (0,3.0) {Symmetric 与 Asymmetric Quantization};

% symmetric
\node[font=\bfseries] at (-3.3,2.25) {Symmetric};
\draw[axis] (-5.4,1.3) -- (-1.2,1.3);
\foreach \x/\lab in {-5.0/$-a$,-3.3/$0$,-1.6/$a$} {
  \draw[tick] (\x,1.18) -- (\x,1.42);
  \node[font=\scriptsize] at (\x,0.92) {\lab};
}
\foreach \x in {-5.0,-4.55,-4.1,-3.65,-3.3,-2.95,-2.5,-2.05,-1.6} {
  \node[qpt] at (\x,1.3) {};
}

% asymmetric
\node[font=\bfseries] at (3.0,2.25) {Asymmetric};
\draw[axis] (0.8,1.3) -- (5.4,1.3);
\foreach \x/\lab in {1.2/$x_{\min}$,2.1/$0$,5.0/$x_{\max}$} {
  \draw[tick] (\x,1.18) -- (\x,1.42);
  \node[font=\scriptsize] at (\x,0.92) {\lab};
}
\foreach \x in {1.2,1.68,2.16,2.64,3.12,3.60,4.08,4.56,5.04} {
  \node[qpt, fill=green!12] at (\x,1.3) {};
}

\node[align=left, text width=10.5cm] at (0,-0.35) {
  symmetric quantization 以 0 为中心，权重中常用；
  asymmetric quantization 可以更好覆盖非对称激活分布，但 zero point 会增加 kernel 复杂度。
};
```

\end{tikzpicture}
\caption{对称与非对称量化的区间映射}
\label{fig:symmetric-asymmetric-quantization}
\end{figure}

\subsection{per-tensor、per-channel、per-group scale}
\label{subsec:per-tensor-channel-group-scale}

Scale 的粒度决定了量化误差和 metadata 开销。

\textbf{Per-tensor quantization} 对整个 tensor 使用一个 scale：

[
s_{\mathrm{tensor}}.
]

优点是 metadata 少、实现简单；缺点是如果 tensor 不同通道分布差异大，一个 scale 难以兼顾所有通道。

\textbf{Per-channel quantization} 对每个输出通道或输入通道使用独立 scale。例如 Linear 权重：

[
\mat{W}\in\mathbb{R}^{H_{\mathrm{out}}\times H_{\mathrm{in}}}.
]

若按 output channel 量化：

[
s_i,\quad i=0,\ldots,H_{\mathrm{out}}-1.
]

第 $i$ 行权重使用 scale $s_i$。这种方式通常显著降低误差。

\textbf{Per-group quantization} 将通道或列划分为小组，每组一个 scale。例如 group size 为 $G$，则每 $G$ 个元素共享一个 scale：

[
s_{i,g}.
]

INT4 LLM weight-only quantization 中 per-group scale 很常见。它在精度和 metadata 之间折中：group 越小，量化更精细，但 scale metadata 越多，kernel 更复杂；group 越大，metadata 少，但误差增大。

[
\text{per-tensor}
\rightarrow
\text{per-channel}
\rightarrow
\text{per-group}
]

不是简单单调关系。实际选择取决于 weight layout、kernel 支持和精度目标。

\begin{table}[H]
\centering
\small
\caption{Scale 粒度对比}
\label{tab:scale-granularity}
\begin{tabular}{p{0.22\textwidth}p{0.30\textwidth}p{0.36\textwidth}}
\toprule
\textbf{粒度} & \textbf{优点} & \textbf{缺点} \
\midrule
Per-tensor
& metadata 最少，实现简单
& outlier 影响全 tensor，误差可能较大。 \
\midrule
Per-channel
& 每个通道适配自身范围，精度较好
& scale 数量增加，kernel 需支持通道 scale。 \
\midrule
Per-group
& INT4 常用，精度和 metadata 折中
& group size 需要调参，解包和 scale 应用更复杂。 \
\bottomrule
\end{tabular}
\end{table}

\subsection{dequantization}
\label{subsec:dequantization}

低精度权重可以有两种计算方式。

第一种是\textbf{反量化后计算}：

[
\hat{\mat{W}}=s(\mat{Q}_W-z),
]
[
\mat{Y}=\mat{X}\hat{\mat{W}}.
]

这种方式实现简单，但如果先把整块权重反量化到 FP16/BF16 再 GEMM，会失去显存带宽收益，并增加临时显存。

第二种是\textbf{融合反量化与 GEMM}。Kernel 从 HBM 读取 packed INT4/INT8 权重，在寄存器或片上存储中解包并乘 scale，然后参与矩阵乘法：

[
\mat{Y}
=======

\mat{X}
\left[
s(\mat{Q}_W-z)
\right].
]

高性能 weight-only 推理必须使用这种方式。也就是说，量化模型的速度不仅取决于量化算法，还取决于是否有高质量 kernel。

\begin{figure}[H]
\centering
\begin{tikzpicture}[
font=\small,
box/.style={draw, rounded corners=3pt, minimum width=2.05cm, minimum height=0.7cm, align=center},
arrow/.style={-{Latex[length=2.2mm]}, thick}
]

```
\node[font=\bfseries] at (0,3.1) {反量化路径：离线展开 vs fused dequant GEMM};

\node[box, fill=blue!8] (q1) at (-5.0,1.55) {INT4 packed\\weights};
\node[box, fill=orange!14] (dq1) at (-2.6,1.55) {dequant to\\FP16 tensor};
\node[box, fill=green!10] (gemm1) at (-0.2,1.55) {FP16 GEMM};
\draw[arrow] (q1) -- (dq1);
\draw[arrow] (dq1) -- (gemm1);

\node[box, fill=blue!8] (q2) at (-5.0,0.2) {INT4 packed\\weights};
\node[box, fill=purple!10, minimum width=3.0cm] (fused) at (-1.9,0.2) {fused unpack + scale + GEMM};
\node[box, fill=green!10] (out) at (1.3,0.2) {output};
\draw[arrow] (q2) -- (fused);
\draw[arrow] (fused) -- (out);

\node[align=left, text width=10.5cm] at (0,-1.35) {
  若先反量化完整 FP16 权重，再执行普通 GEMM，会失去许多带宽收益。
  高性能量化推理通常把 unpack、scale 和 GEMM 融合到一个 kernel 中。
};
```

\end{tikzpicture}
\caption{量化权重的两种计算路径}
\label{fig:dequantization-paths}
\end{figure}

\section{PTQ 与 QAT}
\label{sec:ptq-and-qat}

量化方法可以从训练介入程度分为两大类：

\begin{itemize}
\item \textbf{PTQ, Post-Training Quantization}：训练完成后量化，不重新训练或只做少量校准；
\item \textbf{QAT, Quantization-Aware Training}：训练或微调时模拟量化误差，让模型适应低精度。
\end{itemize}

LLM 推理部署中，PTQ 非常流行，因为重新训练大模型成本很高。GPTQ、AWQ、SmoothQuant 都属于后训练量化或校准型方法。QAT 精度潜力更高，但成本也更高，且对训练数据、训练 recipe 和工程流程要求更强。

\subsection{Post-Training Quantization}
\label{subsec:post-training-quantization}

PTQ 的基本流程是：

\begin{enumerate}
\item 加载训练好的高精度模型；
\item 准备少量校准数据；
\item 收集权重和激活统计；
\item 选择 scale、zero point、group size、clipping 策略；
\item 量化权重或激活；
\item 评估 perplexity、benchmark、下游任务和生成质量；
\item 导出量化模型和推理 kernel 所需 metadata。
\end{enumerate}

PTQ 的优点是：

\begin{itemize}
\item 不需要重新训练；
\item 成本低；
\item 部署速度快；
\item 适合已有开源模型和内部 checkpoint；
\item 可以对不同硬件生成不同量化格式。
\end{itemize}

缺点是：

\begin{itemize}
\item 低 bit 下精度可能明显下降；
\item 对校准数据敏感；
\item outlier 和敏感层难处理；
\item 不同模型架构量化难度不同；
\item 需要专门算法降低误差。
\end{itemize}

\subsection{calibration dataset}
\label{subsec:calibration-dataset}

校准数据用于估计激活范围、outlier、通道重要性和 scale。它不一定需要很大，但必须代表模型实际使用分布。如果部署场景是代码生成，用普通百科文本校准可能不理想；如果部署场景是长文档问答，短指令数据可能无法覆盖长上下文激活分布。

校准数据需要考虑：

\begin{itemize}
\item 领域是否匹配；
\item prompt 长度是否匹配；
\item 是否包含真实 chat template；
\item 是否包含 system prompt；
\item 是否覆盖代码、数学、多语言等高风险输入；
\item 是否覆盖长上下文；
\item 数据量是否足够稳定估计统计量。
\end{itemize}

校准样本数量不是越多越好。太少统计不稳，太多校准耗时高。很多 PTQ 方法使用几百到几千条样本即可得到可用结果，但具体取决于模型和 bit-width。

\subsection{Quantization-Aware Training}
\label{subsec:quantization-aware-training}

QAT 在训练或微调中引入 fake quantization。Forward 时模拟量化：

[
\hat{x}
=======

s\cdot
\operatorname{round}
\left(
\frac{x}{s}
\right),
]

但 backward 时通常使用 straight-through estimator，简称 STE：

[
\frac{\partial \hat{x}}{\partial x}\approx 1.
]

这样模型可以在训练中适应量化误差。QAT 通常比 PTQ 精度更好，尤其在极低 bit、activation quantization 或任务敏感场景中。但 QAT 需要训练资源，也可能需要原始训练数据或高质量微调数据。

QAT 适合：

\begin{itemize}
\item 对质量要求极高的生产模型；
\item 低 bit 下 PTQ 精度不够；
\item 硬件固定，需要深度适配；
\item 模型训练流程可控；
\item 有足够校准/微调数据。
\end{itemize}

\begin{lstlisting}[language=Python,caption={Fake Quantization 的简化 PyTorch 实现},label={lst:fake-quantization-pytorch}]
import torch

class FakeQuantize(torch.autograd.Function):
@staticmethod
def forward(ctx, x, scale, qmin, qmax):
q = torch.round(x / scale)
q = torch.clamp(q, qmin, qmax)
return q * scale

```
@staticmethod
def backward(ctx, grad_output):
    # Straight-through estimator.
    return grad_output, None, None, None
```

def fake_quant(x, scale, qmin=-127, qmax=127):
return FakeQuantize.apply(x, scale, qmin, qmax)
\end{lstlisting}

这个代码只是说明 STE 思想。真实 QAT 需要 observer、scale 更新、per-channel quantization、梯度稳定性、混合精度和训练 recipe 配合。

\subsection{PTQ 与 QAT 的工程选择}
\label{subsec:ptq-vs-qat-engineering-choice}

工程上常见选择是：

[
\text{先 PTQ，若质量不够，再考虑 QAT 或低秩微调补偿。}
]

原因是 PTQ 成本低，可以快速得到多个候选量化配置。例如：

[
\text{W8A16},\quad
\text{W4A16 group size 128},\quad
\text{W8A8 SmoothQuant},\quad
\text{FP8 KV}.
]

如果某个配置满足质量和性能要求，就无需进行昂贵的 QAT。如果 PTQ 在关键任务上损失明显，再考虑 QAT、LoRA-QAT、混合精度保护敏感层等方案。

\begin{table}[H]
\centering
\small
\caption{PTQ 与 QAT 对比}
\label{tab:ptq-vs-qat}
\begin{tabular}{p{0.20\textwidth}p{0.34\textwidth}p{0.34\textwidth}}
\toprule
\textbf{维度} & \textbf{PTQ} & \textbf{QAT} \
\midrule
训练成本
& 低，只需校准或重构
& 高，需要训练或微调 \
\midrule
部署速度
& 快
& 慢 \
\midrule
精度潜力
& 中到高，取决于算法和 bit-width
& 通常更高 \
\midrule
数据需求
& 少量校准数据
& 微调或训练数据 \
\midrule
工程复杂度
& 中等
& 高 \
\midrule
适合场景
& 快速部署、已有模型、多硬件格式
& 极致低 bit、质量要求极高、训练流程可控 \
\bottomrule
\end{tabular}
\end{table}

\section{Weight-only Quantization}
\label{sec:weight-only-quantization}

Weight-only quantization 是 LLM 推理中最常见的量化方式之一。它只量化权重，激活仍保持 FP16/BF16：

[
\mat{Y}
=======

\mat{X}*{\mathrm{fp16}}
\cdot
\operatorname{Dequant}(\mat{W}*{\mathrm{int4/int8}}).
]

常见形式包括：

[
\text{W8A16},
\quad
\text{W4A16}.
]

其中 W 表示 weight bit-width，A 表示 activation bit-width。

\subsection{W8A16 / W4A16}
\label{subsec:w8a16-w4a16}

W8A16 使用 INT8 权重和 FP16/BF16 激活。它通常精度较稳，压缩率约 2 倍，kernel 实现也相对容易。

W4A16 使用 INT4 权重和 FP16/BF16 激活。它压缩率约 4 倍，但更容易损失精度，需要更复杂算法和 kernel 支持。W4A16 常用于单机部署更大模型，或在显存固定时提高并发和上下文长度。

在 weight-only quantization 中，矩阵乘法可以写作：

[
\mat{Y}
=======

\mat{X}
\left(
\mat{S}\odot
\mat{Q}_W
\right),
]

其中 $\mat{Q}_W$ 是低精度整数权重，$\mat{S}$ 是 scale。根据 scale 粒度不同，$\mat{S}$ 可以按 tensor、channel 或 group broadcast。

\subsection{group-wise quantization}
\label{subsec:group-wise-weight-quantization}

INT4 常用 group-wise quantization。设权重矩阵某一行被划分为多个 group，每个 group 大小为 $G$。第 $g$ 个 group 的权重为：

[
\vect{w}_g\in\mathbb{R}^{G}.
]

Symmetric INT4 量化可取：

[
s_g=
\frac{\max_i |w_{g,i}|}{7}.
]

量化：

[
q_{g,i}
=======

\operatorname{clip}
\left(
\left\lfloor
\frac{w_{g,i}}{s_g}
\right\rceil,
-8,
7
\right).
]

反量化：

[
\hat{w}*{g,i}=s_gq*{g,i}.
]

Group size 常见选择如 32、64、128。Group 越小，scale 更精细，精度更好；但 scale metadata 更多，kernel 读取 scale 更频繁。Group 越大，metadata 少，但 outlier 更容易拉大量化范围。

\begin{figure}[H]
\centering
\begin{tikzpicture}[
font=\small,
cell/.style={draw, minimum width=0.38cm, minimum height=0.34cm},
groupbox/.style={draw, rounded corners=2pt, dashed, thick},
arrow/.style={-{Latex[length=2.0mm]}, thick}
]

```
\node[font=\bfseries] at (0,3.0) {Group-wise INT4 Weight Quantization};

\foreach \i in {0,...,23} {
  \node[cell, fill=blue!10] at ({-4.6+0.38*\i},1.45) {};
}

\draw[groupbox] (-4.82,1.17) rectangle (-1.62,1.72);
\draw[groupbox] (-1.54,1.17) rectangle (1.66,1.72);
\draw[groupbox] (1.74,1.17) rectangle (4.94,1.72);

\node[font=\scriptsize] at (-3.22,0.9) {scale $s_0$};
\node[font=\scriptsize] at (0.06,0.9) {scale $s_1$};
\node[font=\scriptsize] at (3.34,0.9) {scale $s_2$};

\node[align=left, text width=10.5cm] at (0,-0.65) {
  每个 group 使用独立 scale。group size 越小，量化误差通常越低；
  但 scale metadata 和 kernel 复杂度也更高。
};
```

\end{tikzpicture}
\caption{按 group 量化权重并保存 scale}
\label{fig:group-wise-weight-quantization}
\end{figure}

\subsection{packing 与 kernel}
\label{subsec:int4-packing-and-kernel}

INT4 权重需要 packing。一个 8-bit byte 可以存两个 4-bit 权重：

[
\text{byte}
===========

(q_0\ &\ 0xF)
|
((q_1\ &\ 0xF)\ll 4).
]

Packing 的目的包括：

\begin{itemize}
\item 降低存储大小；
\item 降低 HBM 读取字节数；
\item 对齐 kernel 的向量化加载；
\item 适配 Tensor Core 或自定义 MMA 路径。
\end{itemize}

下面是教学版 INT4 packing 代码。

\begin{lstlisting}[language=Python,caption={教学版 INT4 packing 与 unpacking},label={lst:int4-pack-unpack}]
import torch

def pack_int4(q):
"""
q: int tensor with values in [0, 15], shape [..., N]
N must be even.
return: uint8 tensor shape [..., N // 2]
"""
assert q.shape[-1] % 2 == 0
q = q.to(torch.uint8)

```
low = q[..., 0::2] & 0x0F
high = (q[..., 1::2] & 0x0F) << 4

return low | high
```

def unpack_int4(packed):
"""
packed: uint8 tensor shape [..., N_packed]
return: uint8 tensor with values in [0, 15], shape [..., 2 * N_packed]
"""
low = packed & 0x0F
high = (packed >> 4) & 0x0F

```
out = torch.empty(
    *packed.shape[:-1],
    packed.shape[-1] * 2,
    dtype=torch.uint8,
    device=packed.device,
)
out[..., 0::2] = low
out[..., 1::2] = high
return out
```

\end{lstlisting}

真实 kernel 不会先 unpack 整个权重矩阵，而是在 GEMM 内部按 tile 解包、转换、乘 scale，并与 activation 做乘加。否则 INT4 的带宽优势会被临时展开抵消。

\subsection{accuracy drop 的来源}
\label{subsec:weight-only-accuracy-drop-sources}

Weight-only quantization 的精度下降主要来自：

\begin{itemize}
\item 权重分布 outlier；
\item 通道重要性差异；
\item 量化误差在层间累积；
\item attention projection 或 FFN 某些矩阵更敏感；
\item 最前/最后几层更敏感；
\item embedding/lm head 更敏感；
\item group size 过大；
\item calibration 数据不匹配；
\item 低 bit 下 rounding error 太大。
\end{itemize}

一个常见工程策略是混合精度保护敏感部分。例如：

\begin{itemize}
\item embedding 保持 FP16；
\item lm head 保持 FP16；
\item 第一层和最后一层保持 FP16/INT8；
\item 某些 outlier channels 使用更高精度；
\item attention output 或 FFN down projection 使用更保守量化。
\end{itemize}

这类策略会牺牲部分压缩率，但常能显著改善质量。

\section{GPTQ}
\label{sec:gptq}

GPTQ 是 LLM PTQ 中非常有代表性的方法。它属于 weight-only post-training quantization，核心思想是利用近似二阶信息，在逐列量化权重时补偿误差，使量化后的层输出尽量接近原始层输出。

\subsection{二阶误差直觉}
\label{subsec:gptq-second-order-intuition}

对某个 Linear 层：

[
\mat{Y}=\mat{X}\mat{W}.
]

量化后权重为：

[
\hat{\mat{W}}=\mat{W}+\Delta\mat{W}.
]

输出误差为：

[
\Delta\mat{Y}
=============

\mat{X}\Delta\mat{W}.
]

我们希望最小化：

[
|\mat{X}\mat{W}-\mat{X}\hat{\mat{W}}|_2^2
=========================================

|\mat{X}\Delta\mat{W}|_2^2.
]

展开可得误差与输入协方差有关：

[
|\mat{X}\Delta\mat{W}|_2^2
==========================

\operatorname{tr}
\left(
\Delta\mat{W}^{\mathsf{T}}
\mat{X}^{\mathsf{T}}\mat{X}
\Delta\mat{W}
\right).
]

令：

[
\mat{H}\approx \mat{X}^{\mathsf{T}}\mat{X}.
]

则量化误差可以用 Hessian-like 矩阵 $\mat{H}$ 加权。直觉是：如果某个权重方向对输出影响大，则量化误差应该更小；如果某个方向对输出影响小，则可以承受更大误差。

\subsection{逐列量化与误差补偿}
\label{subsec:gptq-sequential-quantization}

GPTQ 的一个关键思想是逐列或分块量化权重，并在量化某一列后，把误差补偿到尚未量化的列上。

设当前要量化权重列 $\vect{w}_i$，量化后为 $\hat{\vect{w}}_i$，误差为：

[
\vect{e}_i=\vect{w}_i-\hat{\vect{w}}_i.
]

如果直接忽略误差，后续输出会偏离。GPTQ 利用 $\mat{H}^{-1}$ 的信息，把该误差对后续列的影响进行更新。抽象地：

[
\mat{W}*{j}
\leftarrow
\mat{W}*{j}
-----------

\alpha_{ij}\vect{e}_i,
\quad j>i.
]

其中 $\alpha_{ij}$ 来自 Hessian inverse 的相关项。这样后续列在量化前已经吸收部分误差补偿。

这个过程类似“带二阶信息的贪心量化”。它不是简单地对每个 group 独立 round，而是考虑输入分布和权重列之间的耦合。

\begin{figure}[H]
\centering
\begin{tikzpicture}[
font=\small,
col/.style={draw, rounded corners=2pt, minimum width=0.55cm, minimum height=1.6cm, align=center},
arrow/.style={-{Latex[length=2.1mm]}, thick}
]

```
\node[font=\bfseries] at (0,3.0) {GPTQ：逐列量化并向后续列传播误差补偿};

\foreach \i/\x/\colr in {0/-3.6/green!12,1/-2.9/blue!10,2/-2.2/blue!10,3/-1.5/blue!10,4/-0.8/blue!10} {
  \node[col, fill=\colr] (c\i) at (\x,1.25) {$w_{\i}$};
}

\node[draw, rounded corners=3pt, fill=orange!14, minimum width=1.4cm, minimum height=0.6cm, align=center] (q) at (-3.6,-0.1)
  {quantize\\$w_0$};

\draw[arrow] (c0) -- (q);

\foreach \i in {1,2,3,4} {
  \draw[arrow, red!70] (q) .. controls (-3.0,-0.85) and (-2.2,-0.85) .. (c\i.south);
}

\node[align=left, text width=10.5cm] at (0,-1.5) {
  GPTQ 不是独立量化每一列，而是在量化当前列后，将误差按近似二阶信息补偿到后续未量化列。
  这样可以降低层输出重构误差。
};
```

\end{tikzpicture}
\caption{GPTQ 的逐列量化与误差补偿直觉}
\label{fig:gptq-sequential-error-compensation}
\end{figure}

\subsection{Hessian 近似}
\label{subsec:gptq-hessian-approximation}

GPTQ 不会计算完整训练 loss 的真实 Hessian，而是对每层局部重构误差使用近似：

[
\mat{H}=\mat{X}^{\mathsf{T}}\mat{X}.
]

这里 $\mat{X}$ 是校准数据经过该层时的输入 activation。实际实现中会加入 damping，避免矩阵病态：

[
\mat{H}_{\mathrm{damped}}
=========================

\mat{H}
+
\lambda \operatorname{diag}(\mat{H}).
]

然后使用 Cholesky 分解或其他数值方法得到所需的逆相关信息。由于层维度很大，GPTQ 通常采用 blockwise 处理，避免一次操作完整巨大矩阵。

\subsection{工程流程}
\label{subsec:gptq-engineering-flow}

GPTQ 的工程流程可以概括为：

\begin{enumerate}
\item 准备校准样本；
\item 对模型逐层 forward，收集每层输入 activation；
\item 对每个 Linear 层构造 $\mat{H}\approx \mat{X}^{\mathsf{T}}\mat{X}$；
\item 对权重按 block/group 逐步量化；
\item 使用 Hessian inverse 信息进行误差补偿；
\item 导出 INT4/INT3/INT8 权重、scale、zero point 和 packing layout；
\item 使用支持 GPTQ 格式的推理 kernel 执行。
\end{enumerate}

GPTQ 的优势是低 bit 下精度较好，尤其适合 weight-only INT4。代价是量化过程比简单 min-max PTQ 更复杂，需要校准数据和矩阵分解，量化耗时更长。

\section{AWQ}
\label{sec:awq}

AWQ，Activation-aware Weight Quantization，是另一类重要的 LLM weight-only PTQ 方法。它的核心观察是：权重并非同等重要，少量 salient weight channels 对模型输出影响很大；识别这些重要通道时，应该关注 activation 分布，而不只是权重本身。

\subsection{activation-aware 的动机}
\label{subsec:awq-activation-aware-motivation}

对 Linear 层：

[
\mat{Y}=\mat{X}\mat{W}.
]

某个权重通道的量化误差对输出影响，不仅取决于权重误差本身，还取决于输入 activation 的大小。如果某个输入通道的 activation 经常很大，那么对应权重列的误差会被放大。

设第 $j$ 个输入通道 activation 平均幅度较大：

[
\mathbb{E}[|X_j|]\ \text{large}.
]

那么权重列 $\mat{W}_{j,:}$ 的量化误差会对输出产生更大影响：

[
\Delta\mat{Y}
=============

\sum_j
\mat{X}*j
\Delta\mat{W}*{j,:}.
]

因此，AWQ 使用 activation 统计来判断哪些 channels 更重要，并对这些 channels 做保护。

\subsection{salient channels}
\label{subsec:awq-salient-channels}

AWQ 的思想可以简化为：

[
\text{找出 activation 大、对输出敏感的 channels}
\rightarrow
\text{降低这些 channels 的量化误差}.
]

一种直接做法是混合精度：把重要 channels 保持高精度。但混合精度可能不硬件友好，因为 kernel 需要同时处理 INT4 和 FP16 权重，导致实现复杂。

AWQ 更偏向硬件友好方案：通过缩放变换保护 salient channels，同时保持整体权重仍可统一低 bit 量化。

\subsection{scale 变换}
\label{subsec:awq-scale-transformation}

考虑 Linear：

[
\mat{Y}=\mat{X}\mat{W}.
]

引入对角缩放矩阵 $\mat{S}$：

[
\mat{Y}
=======

(\mat{X}\mat{S}^{-1})
(\mat{S}\mat{W}).
]

这是数学等价变换：

[
\mat{X}\mat{W}
==============

\mat{X}\mat{S}^{-1}\mat{S}\mat{W}.
]

AWQ 的直觉是，适当放大某些重要 weight channels，使其在量化时相对误差更小，同时在 activation 侧用反向缩放保持等价。实际中，需要搜索或确定合适的 scale，使量化后误差最小。

这个方法的优点是：最终仍然可以使用统一的 low-bit weight-only kernel，而不必在 kernel 内混合处理不同精度的权重。

\begin{figure}[H]
\centering
\begin{tikzpicture}[
font=\small,
box/.style={draw, rounded corners=3pt, minimum width=1.8cm, minimum height=0.65cm, align=center},
arrow/.style={-{Latex[length=2.1mm]}, thick}
]

```
\node[font=\bfseries] at (0,3.05) {AWQ：通过等价缩放保护重要通道};

\node[box, fill=blue!8] (x) at (-5.0,1.3) {$X$};
\node[box, fill=orange!14] (sinv) at (-2.9,1.3) {$S^{-1}$\\on activation};
\node[box, fill=green!10] (sweight) at (-0.6,1.3) {$S W$\\scaled weight};
\node[box, fill=purple!10] (quant) at (1.7,1.3) {INT4\\quantize};
\node[box, fill=red!8] (y) at (4.2,1.3) {$Y$};

\draw[arrow] (x) -- (sinv);
\draw[arrow] (sinv) -- (sweight);
\draw[arrow] (sweight) -- (quant);
\draw[arrow] (quant) -- (y);

\node[align=left, text width=10.5cm] at (0,-0.65) {
  AWQ 使用 activation 统计识别重要通道，并通过等价缩放降低这些通道的量化误差。
  最终仍保持硬件友好的 weight-only quantization。
};
```

\end{tikzpicture}
\caption{AWQ 中 activation-aware scaling 的直觉}
\label{fig:awq-scale-transformation}
\end{figure}

\subsection{AWQ 的适用场景}
\label{subsec:awq-use-cases}

AWQ 适合：

\begin{itemize}
\item INT4 weight-only 推理；
\item 希望校准成本较低；
\item 希望保持 kernel 硬件友好；
\item 模型是 instruction-tuned LLM 或多模态 LLM；
\item 部署在边缘设备或显存紧张环境；
\item 不希望做复杂二阶重构。
\end{itemize}

与 GPTQ 相比，AWQ 更强调 activation-aware channel protection 和硬件友好性；GPTQ 更强调二阶误差补偿和逐列重构。实际中二者都常作为 INT4 量化候选，需要在目标模型和任务上评估。

\section{SmoothQuant 与 W8A8}
\label{sec:smoothquant-w8a8}

Weight-only quantization 只量化权重，激活仍是 FP16/BF16。若想进一步利用 INT8 GEMM，需要同时量化权重和激活，即 W8A8。SmoothQuant 是 LLM W8A8 PTQ 中代表性方法之一。

\subsection{activation outlier 问题}
\label{subsec:activation-outlier-problem}

LLM 激活中常存在 outlier channels。某些通道的 activation 幅度远大于其他通道。如果直接对 activation 做 per-tensor INT8 量化，scale 会被 outlier 拉大，导致大多数普通值分辨率不足。

设 activation 向量某些通道非常大：

[
|x_j|\gg |x_k|.
]

Per-tensor scale 必须覆盖最大值：

[
s=\frac{\max_j |x_j|}{127}.
]

若 $\max_j |x_j|$ 很大，普通通道的值会被量化到很粗的整数网格上，误差增大。

权重相比激活更容易量化，因为权重是固定的，可以按通道或 group 使用 scale；activation 是动态的，范围随输入变化，outlier 更难处理。

\subsection{将量化难度从 activation 转移到 weight}
\label{subsec:smoothquant-migrate-difficulty}

SmoothQuant 的核心思想是使用等价变换，将 activation outlier 的难度迁移到 weight 上。

对于 Linear：

[
\mat{Y}=\mat{X}\mat{W}.
]

引入对角矩阵 $\mat{S}$：

[
\mat{Y}
=======

(\mat{X}\mat{S}^{-1})
(\mat{S}\mat{W}).
]

如果某些 activation channel 很大，可以选择较大的 scale $s_j$，让：

[
\mat{X}'_j=\frac{\mat{X}_j}{s_j}
]

变得更平滑。对应权重：

[
\mat{W}'_j=s_j\mat{W}_j.
]

这样 activation 更容易 INT8 量化，而权重量化难度增加。由于权重是静态的，更容易通过 per-channel scale 处理。因此整体 W8A8 可行性提高。

SmoothQuant 通常用参数 $\alpha$ 控制平滑程度。若 activation 通道幅度为 $a_j$，weight 通道幅度为 $w_j$，scale 可设计为某种函数：

[
s_j
===

\frac{a_j^{\alpha}}
{w_j^{1-\alpha}}.
]

直觉上，$\alpha$ 越大，越多地平滑 activation；$\alpha$ 越小，越少改变权重。

\begin{figure}[H]
\centering
\begin{tikzpicture}[
font=\small,
box/.style={draw, rounded corners=3pt, minimum width=1.95cm, minimum height=0.7cm, align=center},
arrow/.style={-{Latex[length=2.1mm]}, thick}
]

```
\node[font=\bfseries] at (0,3.1) {SmoothQuant：平滑 Activation Outlier};

\node[box, fill=blue!8] (x) at (-5.1,1.45) {activation\\with outliers};
\node[box, fill=orange!14] (scale) at (-2.7,1.45) {$X'=XS^{-1}$\\smooth};
\node[box, fill=green!10] (w) at (-0.2,1.45) {$W'=SW$\\adjust weight};
\node[box, fill=purple!10] (w8a8) at (2.3,1.45) {W8A8\\GEMM};
\node[box, fill=red!8] (y) at (4.8,1.45) {$Y$};

\draw[arrow] (x) -- (scale);
\draw[arrow] (scale) -- (w8a8);
\draw[arrow] (w) -- (w8a8);
\draw[arrow] (w8a8) -- (y);

\node[align=left, text width=10.5cm] at (0,-0.65) {
  SmoothQuant 通过等价缩放把 activation outlier 平滑掉，
  将部分量化难度转移到静态权重上，从而支持 W8A8 PTQ。
};
```

\end{tikzpicture}
\caption{SmoothQuant 的 activation 平滑思想}
\label{fig:smoothquant-activation-smoothing}
\end{figure}

\subsection{W8A8 GEMM}
\label{subsec:w8a8-gemm}

W8A8 的目标是让权重和激活都用 INT8 表示，矩阵乘法使用 INT8 Tensor Core 或等价硬件路径：

[
\mat{Y}_{\mathrm{int32}}
========================

\mat{X}*{\mathrm{int8}}
\mat{W}*{\mathrm{int8}}.
]

累积通常使用 INT32，然后再乘 scale 转回 FP16/BF16 或继续保持低精度：

[
\hat{\mat{Y}}
=============

s_Xs_W
\mat{Y}_{\mathrm{int32}}.
]

W8A8 的潜在收益包括：

\begin{itemize}
\item 权重读取减半；
\item 激活读取减半；
\item 使用 INT8 矩阵乘加硬件；
\item 降低中间 activation 显存。
\end{itemize}

难点包括：

\begin{itemize}
\item activation 动态范围；
\item outlier；
\item scale 管理；
\item layernorm/RMSNorm 后分布变化；
\item residual add 的 scale 对齐；
\item kernel 对动态 scale 的支持；
\item 不同层量化敏感度差异。
\end{itemize}

\section{FP8 与低精度浮点}
\label{sec:fp8-low-precision-floating-point}

INT 量化使用固定线性 scale 映射到整数。FP8 则使用 8 bit 浮点格式，包含 sign、exponent 和 mantissa。相比 INT8，FP8 可以覆盖更大动态范围，适合训练和推理中的激活、权重或梯度。

\subsection{E4M3 与 E5M2}
\label{subsec:e4m3-e5m2}

常见 FP8 格式包括：

\begin{itemize}
\item \textbf{E4M3}：4 bit exponent，3 bit mantissa；
\item \textbf{E5M2}：5 bit exponent，2 bit mantissa。
\end{itemize}

E4M3 尾数更多，精度更高，但动态范围较小。E5M2 指数更多，动态范围更大，但精度较低。不同张量可能适合不同格式。例如 forward activation 和 weight 可能更偏 E4M3，gradient 可能需要更大动态范围。

\begin{figure}[H]
\centering
\begin{tikzpicture}[
font=\small,
bit/.style={draw, minimum width=0.5cm, minimum height=0.45cm, align=center}
]

```
\node[font=\bfseries] at (0,3.0) {FP8 E4M3 与 E5M2 bit layout};

\node[anchor=east] at (-4.2,1.55) {E4M3};
\node[bit, fill=red!15] at (-3.6,1.55) {S};
\foreach \i/\x in {0/-3.05,1/-2.5,2/-1.95,3/-1.4} {
  \node[bit, fill=blue!12] at (\x,1.55) {E};
}
\foreach \i/\x in {0/-0.85,1/-0.3,2/0.25} {
  \node[bit, fill=green!12] at (\x,1.55) {M};
}

\node[anchor=east] at (-4.2,0.65) {E5M2};
\node[bit, fill=red!15] at (-3.6,0.65) {S};
\foreach \i/\x in {0/-3.05,1/-2.5,2/-1.95,3/-1.4,4/-0.85} {
  \node[bit, fill=blue!12] at (\x,0.65) {E};
}
\foreach \i/\x in {0/-0.3,1/0.25} {
  \node[bit, fill=green!12] at (\x,0.65) {M};
}

\node[align=left, text width=10.4cm] at (0,-0.75) {
  E4M3 提供更多 mantissa 精度，E5M2 提供更大 exponent 动态范围。
  FP8 通常需要配合 scale、amax history 和硬件支持使用。
};
```

\end{tikzpicture}
\caption{FP8 两种常见格式的 bit 分配}
\label{fig:fp8-e4m3-e5m2}
\end{figure}

\subsection{scale 与 amax}
\label{subsec:fp8-scale-amax}

FP8 表示范围有限，需要 scale 管理。常见做法是记录 tensor 的绝对最大值：

[
a_{\max}=\max_i |x_i|.
]

根据 FP8 格式最大可表示值 $f_{\max}$，选择 scale：

[
s=
\frac{a_{\max}}{f_{\max}}.
]

量化时：

[
x_{\mathrm{fp8}}
================

\operatorname{cast}_{\mathrm{fp8}}
\left(
\frac{x}{s}
\right).
]

反量化时：

[
\hat{x}=s\cdot x_{\mathrm{fp8}}.
]

为了避免 scale 随 batch 抖动，系统可能维护 amax history，并用滑动窗口或 delayed scaling 更新 scale。这样可以减少异常 batch 导致的 scale 波动。

FP8 的工程难点在于：每个 tensor 或每类 tensor 的 scale 如何选择、何时更新、是否 per-tensor 或 per-channel、是否 delayed、如何与分布式并行和 kernel fusion 配合。

\subsection{FP8 训练与推理}
\label{subsec:fp8-training-inference}

FP8 可用于训练和推理。训练中，部分 GEMM 输入可以使用 FP8，累积通常使用更高精度，master weights 仍可能保留 BF16/FP32。推理中，权重和激活可使用 FP8，以降低带宽和提高吞吐。

FP8 相比 INT8 的优势是动态范围更自然；相比 BF16 的优势是字节数更低。它更依赖硬件支持，例如特定 GPU 架构的 FP8 Tensor Core，以及软件库对 scale 管理和 kernel 的支持。

FP8 与 INT8/INT4 的关系不是谁完全替代谁：

\begin{itemize}
\item INT4 weight-only 适合极致压缩权重；
\item INT8 W8A8 适合成熟整数硬件路径；
\item FP8 适合有硬件支持、希望兼顾动态范围和低字节数的训练/推理；
\item KV Cache 可用 FP8 或 INT8，取决于质量和 kernel 支持。
\end{itemize}

\section{KV Cache 量化}
\label{sec:kv-cache-quantization}

前面几章已经看到，KV Cache 是长上下文和高并发 serving 的关键显存瓶颈。KV Cache 量化的目标是降低：

[
M_{\mathrm{KV}}
===============

2LBSn_{kv}d_hs_{\mathrm{dtype}}.
]

如果把 BF16 KV Cache 换成 FP8 或 INT8，$s_{\mathrm{dtype}}$ 从 2 bytes 降为 1 byte，理论上 KV Cache 显存减半。若使用 INT4，理论上可进一步减半，但质量风险和 kernel 复杂度更高。

\subsection{KV Cache 为什么适合量化}
\label{subsec:why-kv-cache-quantization}

KV Cache 量化适合以下场景：

\begin{itemize}
\item 长上下文；
\item 高并发；
\item batch decode；
\item GQA/MQA 后仍然 cache 很大；
\item 显存主要被 KV Cache 占用；
\item decode memory bandwidth 成为瓶颈。
\end{itemize}

在 decode 中，每一步都读取历史 K/V。如果 K/V 字节数减半，HBM 读取压力也可能下降。尤其对于 memory-bound decode attention，KV 量化有潜在速度收益。

\subsection{K 与 V 的敏感性}
\label{subsec:k-v-sensitivity}

K 和 V 的量化误差影响不同。K 参与 attention score：

[
qk^{\mathsf{T}}.
]

K 的误差会改变 attention 分布：

[
\softmax(q(K+\Delta K)^{\mathsf{T}}).
]

V 的误差会改变被加权汇聚的内容：

[
P(V+\Delta V).
]

K 的误差可能通过 softmax 放大，因为 score 的相对差异会影响注意力权重；V 的误差则更像输出内容误差。实际系统可能对 K 和 V 使用不同 scale、不同格式，甚至只量化 V 或对 K 更保守。

\subsection{per-token 与 per-channel KV scale}
\label{subsec:kv-per-token-per-channel-scale}

KV Cache 量化需要 scale。常见粒度包括：

\begin{itemize}
\item per-tensor：整个 cache 一个 scale，简单但误差大；
\item per-layer：每层一个 scale；
\item per-head：每个 KV head 一个 scale；
\item per-token：每个 token 的 K/V 向量一个 scale；
\item per-channel：每个 head dimension 一个 scale；
\item per-block：PagedAttention block 内共享 scale。
\end{itemize}

Per-token scale 可以适应不同 token 的幅度变化，但 metadata 较多；per-channel scale 可以保护某些维度，但 kernel 读取 scale 更复杂；per-block scale 与 paged KV blocks 很自然，metadata 与 block table 可结合。

KV Cache 量化的工程设计需要同时考虑：

\begin{itemize}
\item scale metadata 显存；
\item attention kernel 读取 scale 的开销；
\item PagedAttention block layout；
\item TP 下 KV heads 切分；
\item prefix cache 共享 blocks；
\item copy-on-write；
\item 长上下文质量。
\end{itemize}

\begin{figure}[H]
\centering
\begin{tikzpicture}[
font=\small,
block/.style={draw, rounded corners=2pt, minimum width=0.7cm, minimum height=0.45cm, align=center},
scale/.style={draw, rounded corners=2pt, minimum width=0.9cm, minimum height=0.35cm, align=center, fill=orange!14}
]

```
\node[font=\bfseries] at (0,3.0) {Paged KV Cache 中的 per-block scale};

\foreach \i/\x in {0/-3.6,1/-2.7,2/-1.8,3/-0.9,4/0.0} {
  \node[block, fill=blue!10] (b\i) at (\x,1.35) {$B_{\i}$};
  \node[scale] at (\x,0.72) {$s_{\i}$};
}

\foreach \i/\x in {5/1.2,6/2.1,7/3.0} {
  \node[block, fill=green!10] (b\i) at (\x,1.35) {$B_{\i}$};
  \node[scale] at (\x,0.72) {$s_{\i}$};
}

\node[align=left, text width=10.5cm] at (0,-0.75) {
  KV Cache 量化可以为每个 physical block 保存 scale。
  这种方式与 PagedAttention 的 block 管理自然结合，但 attention kernel 必须读取并应用对应 scale。
};
```

\end{tikzpicture}
\caption{KV Cache per-block scale 与 paged blocks}
\label{fig:kv-cache-per-block-scale}
\end{figure}

\section{量化对 Serving Infra 的影响}
\label{sec:quantization-impact-on-serving-infra}

量化不是只改变模型文件格式。它会影响模型加载、GPU kernel、调度器、KV Cache、并行推理、监控和质量评估。

\subsection{模型加载与格式}
\label{subsec:quantized-model-loading-format}

量化模型通常包含：

\begin{itemize}
\item packed low-bit weights；
\item scale tensors；
\item zero point tensors；
\item group size metadata；
\item quantization method metadata；
\item layer-wise configuration；
\item kernel expected layout；
\item optional outlier/salient channel metadata。
\end{itemize}

Serving engine 加载模型时必须确保：

\begin{itemize}
\item 权重 layout 与 kernel 匹配；
\item TP 切分与 packed format 兼容；
\item scale 切分与权重 shard 对齐；
\item checkpoint 转换没有破坏 group 边界；
\item lm head、embedding 等特殊层处理正确；
\item adapter 或 LoRA 与量化权重兼容。
\end{itemize}

一个常见错误是：量化权重在单 GPU 上能跑，但切到 TP 后 group scale 对齐出错，导致输出异常。原因是 TP shard 边界可能切开 group，或 scale metadata 没有同步切分。因此，量化格式必须从一开始就考虑并行推理。

\subsection{kernel 支持}
\label{subsec:quantized-kernel-support}

量化收益高度依赖 kernel。一个量化配置在纸面上显存很省，但如果没有高性能 kernel，可能速度很慢。

Kernel 需要支持：

\begin{itemize}
\item INT4/INT8 unpack；
\item per-group/per-channel scale；
\item zero point；
\item fused dequant GEMM；
\item Tensor Parallel shard；
\item activation dtype；
\item batch GEMM shape；
\item CUDA graph 或动态 shape；
\item MoE grouped GEMM；
\item LoRA adapter 合并或旁路。
\end{itemize}

如果推理框架 fallback 到普通 FP16 GEMM，可能会先反量化权重，导致显存和速度都不理想。因此上线前必须通过 profiler 或日志确认量化 kernel 被真正调用。

\subsection{质量评估}
\label{subsec:quantization-quality-evaluation}

量化上线不能只看 perplexity。需要多层评估：

\begin{itemize}
\item 语言建模 perplexity；
\item 通用 benchmark；
\item 指令跟随；
\item 代码能力；
\item 数学推理；
\item 多语言；
\item 长上下文；
\item RAG 场景；
\item 安全策略；
\item 真实线上 prompt 回放；
\item 人工偏好评测。
\end{itemize}

低 bit 量化常见问题包括：

\begin{itemize}
\item 输出重复；
\item 格式遵循下降；
\item 长推理链更容易出错；
\item 代码生成中小语法错误增加；
\item 数学计算稳定性下降；
\item 长上下文引用准确率下降；
\item 多语言弱语言质量下降。
\end{itemize}

因此，量化模型需要按业务场景评估，而不是只看一个平均指标。

\subsection{成本模型}
\label{subsec:quantization-cost-model}

量化的最终目标通常是降低成本。可以定义单位 token 成本：

[
C_{\mathrm{token}}
==================

\frac{C_{\mathrm{GPU/hour}}}
{\mathrm{tokens/hour}}.
]

量化可能通过以下路径降低成本：

\begin{itemize}
\item 用更少 GPU 部署同一模型；
\item 同一 GPU 承载更多并发；
\item 提高 tokens/s；
\item 降低能耗；
\item 留出更多显存给 KV Cache；
\item 支持更长上下文而不扩容。
\end{itemize}

但量化也可能引入额外成本：

\begin{itemize}
\item 质量下降导致重试或人工干预；
\item kernel 不成熟导致延迟变差；
\item 多格式模型管理复杂；
\item 量化校准与评估成本；
\item 特定硬件依赖增加。
\end{itemize}

因此，工程决策应比较端到端成本，而不是只看模型大小。

\section{一个教学版量化实现}
\label{sec:toy-quantization-implementation}

下面给出一个教学版 per-group symmetric INT4 权重量化实现。它用于理解数据结构，不是生产级实现。

\begin{lstlisting}[language=Python,caption={Per-group symmetric INT4 权重量化示意},label={lst:per-group-int4-quantization}]
import torch

def quantize_int4_per_group(weight, group_size=128):
"""
weight: [out_features, in_features], FP16/BF16/FP32
return:
q_packed: packed uint8 INT4 weights
scales: [out_features, num_groups]
"""
out_features, in_features = weight.shape
assert in_features % group_size == 0
assert group_size % 2 == 0

```
num_groups = in_features // group_size

w = weight.float().view(out_features, num_groups, group_size)

max_abs = w.abs().amax(dim=-1).clamp_min(1e-8)
scales = max_abs / 7.0

q = torch.round(w / scales[..., None]).clamp(-8, 7)

# Convert signed INT4 [-8, 7] to unsigned nibble [0, 15].
q_u = (q.to(torch.int16) + 8).to(torch.uint8)

# Pack two int4 values into one byte.
low = q_u[..., 0::2] & 0x0F
high = (q_u[..., 1::2] & 0x0F) << 4
q_packed = (low | high).contiguous()

return q_packed, scales.half()
```

def dequantize_int4_per_group(q_packed, scales, group_size=128):
"""
Slow reference dequantization for correctness tests.
"""
out_features, num_groups, packed_per_group = q_packed.shape
assert packed_per_group * 2 == group_size

```
low = q_packed & 0x0F
high = (q_packed >> 4) & 0x0F

q_u = torch.empty(
    out_features,
    num_groups,
    group_size,
    dtype=torch.uint8,
    device=q_packed.device,
)
q_u[..., 0::2] = low
q_u[..., 1::2] = high

q = q_u.to(torch.int16) - 8
w = q.float() * scales.float()[..., None]

return w.view(out_features, num_groups * group_size)
```

\end{lstlisting}

量化后需要验证重构误差：

[
\frac{|\mat{W}-\hat{\mat{W}}|_F}
{|\mat{W}|_F}.
]

更重要的是验证层输出误差：

[
\frac{|\mat{X}\mat{W}-\mat{X}\hat{\mat{W}}|_F}
{|\mat{X}\mat{W}|_F}.
]

因为模型质量受输出误差影响，而不是只受权重误差影响。

\begin{lstlisting}[language=Python,caption={量化权重的重构误差与层输出误差检查},label={lst:quantization-error-check}]
def check_quantization_error(weight, x, group_size=128):
q_packed, scales = quantize_int4_per_group(weight, group_size)
w_hat = dequantize_int4_per_group(q_packed, scales, group_size)

```
weight_error = torch.norm(weight.float() - w_hat.float())
weight_error = weight_error / torch.norm(weight.float())

y_ref = x.float() @ weight.float().t()
y_hat = x.float() @ w_hat.float().t()

output_error = torch.norm(y_ref - y_hat) / torch.norm(y_ref)

return {
    "relative_weight_error": weight_error.item(),
    "relative_output_error": output_error.item(),
}
```

\end{lstlisting}

生产系统中，不能对每层都用这种慢速反量化路径执行推理。但这些 reference 实现适合单元测试：确认 packing、scale、zero point、group 切分和 TP shard 都正确。

\section{量化方案选型}
\label{sec:quantization-scheme-selection}

不同量化方案适合不同部署目标。

\subsection{显存优先}
\label{subsec:quantization-memory-first}

如果目标是尽可能降低显存，使模型能装进更少 GPU，可优先考虑：

[
\text{W4A16 GPTQ/AWQ}
]

或更低 bit weight-only。适合：

\begin{itemize}
\item 单机部署大模型；
\item 边缘设备；
\item 低并发但显存紧张；
\item 成本敏感场景；
\item 能接受少量质量损失。
\end{itemize}

需要重点验证模型质量和 kernel 是否高效。

\subsection{吞吐优先}
\label{subsec:quantization-throughput-first}

如果目标是提升吞吐，需要看硬件是否支持低精度计算。候选包括：

[
\text{W8A8},
\quad
\text{FP8},
\quad
\text{INT4 fused GEMM}.
]

如果只压缩权重但计算仍受 activation 或 compute 限制，速度收益可能有限。吞吐优先场景应使用 profiler 分析：

\begin{itemize}
\item weight read bandwidth 是否下降；
\item GEMM kernel 是否使用低精度路径；
\item dequantization 是否 fused；
\item decode batch 是否变大；
\item TTFT 和 TPOT 是否改善。
\end{itemize}

\subsection{质量优先}
\label{subsec:quantization-quality-first}

如果质量优先，可考虑：

\begin{itemize}
\item W8A16 而非 W4A16；
\item INT4 中使用更小 group size；
\item 保留敏感层为 FP16；
\item 使用 GPTQ/AWQ 而非 naive min-max；
\item 使用 QAT；
\item 对业务数据做校准；
\item 对长上下文和推理任务做专项评估。
\end{itemize}

质量优先并不意味着不能量化，而是需要保守配置和充分评估。

\subsection{推理框架兼容性}
\label{subsec:inference-framework-compatibility}

量化方案必须与推理框架兼容。选择前需要确认：

\begin{itemize}
\item 框架是否支持该量化格式；
\item 是否支持目标 GPU；
\item 是否支持 TP/PP；
\item 是否支持 PagedAttention；
\item 是否支持 LoRA；
\item 是否支持 speculative decoding；
\item 是否支持 streaming；
\item 是否支持 sharded checkpoint；
\item 是否支持导出和加载 scale metadata；
\item 是否能在生产中稳定运行。
\end{itemize}

一个在离线脚本中可用的量化模型，不一定能直接放进 serving engine。模型格式、kernel、scheduler、KV Cache 和并行策略必须一起验证。

\section{本章小结}
\label{sec:chapter15-summary}

本章系统介绍了大模型推理中的量化技术，从数值表示、PTQ/QAT、weight-only quantization，到 GPTQ、AWQ、SmoothQuant、FP8、KV Cache 量化和 serving infra 影响。

\begin{itemize}
\item \textbf{量化的核心目标} 是用更低 bit-width 表示权重、激活或 KV Cache，从而降低显存占用、减少 HBM 带宽压力，并在硬件支持时提升计算吞吐。
\item \textbf{Weight-only quantization} 是 LLM 推理中最常见的方案之一，典型形式包括 W8A16 和 W4A16。
\item \textbf{Uniform quantization} 使用 scale 和 zero point 将浮点值映射到离散整数；symmetric quantization 简单硬件友好，asymmetric quantization 更适合非对称分布。
\item \textbf{Scale 粒度很关键}：per-tensor 简单但误差大，per-channel 精度较好，per-group 是 INT4 LLM 推理中的常见折中。
\item \textbf{高性能量化推理依赖 fused dequant GEMM}。如果先把低 bit 权重完整反量化为 FP16，再执行普通 GEMM，会损失许多带宽收益。
\item \textbf{PTQ 成本低，适合快速部署已有模型}；QAT 精度潜力更高，但需要训练或微调成本。
\item \textbf{GPTQ 利用近似二阶信息进行逐列量化和误差补偿}，适合低 bit weight-only PTQ。
\item \textbf{AWQ 使用 activation 统计识别重要通道}，通过等价缩放保护 salient channels，同时保持硬件友好的 weight-only quantization。
\item \textbf{SmoothQuant 通过等价变换平滑 activation outliers}，将量化难度从动态 activation 迁移到静态 weight，从而支持 W8A8 PTQ。
\item \textbf{FP8 是低精度浮点格式}，通常需要 scale、amax history 和硬件支持，适合训练和推理中的低精度 GEMM。
\item \textbf{KV Cache 量化可以显著降低长上下文和高并发 serving 的显存压力}，但需要处理 K/V 敏感性、scale 粒度和 paged block layout。
\item \textbf{量化会影响整个 serving stack}：模型格式、kernel、TP/PP 切分、PagedAttention、scheduler、质量评估和成本模型都要一起考虑。
\item \textbf{量化方案选型必须结合目标}：显存优先可选 INT4 weight-only，吞吐优先要看低精度硬件路径，质量优先则需要保守配置、敏感层保护和充分评估。
\end{itemize}

下一章将进入推理服务调度与 Continuous Batching。我们会把前面几章的 KV Cache、PagedAttention、FlashAttention 和量化技术放到一个在线 serving engine 中，分析请求如何排队、prefill 与 decode 如何混部、chunked prefill 如何降低首 token 延迟，以及调度器如何在吞吐、延迟、公平性和显存之间做决策。
